% ============================================================
% Section 5: Results
% DisMech Benchmark: RAG vs. Structured Knowledge
% ============================================================

\section{Results}
\label{sec:results}

\subsection{Experimental Setup}
\label{sec:results-setup}

We evaluated two conditions over ten questions drawn from Categories~3, 5,
and~6 of the query-complexity taxonomy described in
Section~\ref{sec:taxonomy}. Ground truth was computed by a deterministic
Python scan of all 605 non-history DisMech YAML files, independent of any
language model or retrieval system.

\paragraph{Corpus.}
The RAG corpus consists of 8{,}217 unique PubMed abstracts cited across all
DisMech disorder entries, fetched via the NCBI Entrez efetch API, plus
2{,}764 mechanism-description texts and 475 disease-description texts
extracted directly from the YAML files---11{,}456 documents in total. Each
document was embedded with Voyage~AI \texttt{voyage-4-large}
(1{,}024-dimensional vectors) and stored in a Qdrant vector index under the
collection \texttt{dismech\_benchmark}.

\paragraph{RAG condition.}
For each question, the question text was embedded with the same model
(\texttt{input\_type=query}) and the top-20 most similar documents were
retrieved by cosine similarity. The retrieved passages were concatenated as
context and passed to Claude Sonnet~4.6 with the instruction to answer
\emph{using only the provided context} and to state explicitly if the context
was insufficient. No access to the DisMech YAML files or TypeDB graph was
permitted.

\paragraph{Structured condition.}
Each question was answered by a deterministic Python query over the raw YAML
files: regular-expression matching over mechanism-description text for
pathway questions (Q1, Q3); direct field access for categorical and ontology
questions (Q2, Q6); absence detection over structured fields for negative-space
questions (Q4, Q5); and sort-and-limit operations for ranking questions
(Q7--Q10). The exact query result was passed to the same Claude Sonnet~4.6
model, which was instructed to format the structured output as a natural-language
answer. No retrieval step was performed.

\paragraph{Scoring.}
Responses were scored against ground truth using two complementary measures
described in Section~\ref{sec:results-overall}. Both conditions used identical
model, temperature, and maximum token settings; the only difference was the
information source provided as context.

\subsection{Overall Accuracy}
\label{sec:results-overall}

Table~\ref{tab:accuracy} reports accuracy for both conditions across all ten
benchmark questions. We report two measures: \emph{binary accuracy} (did the
condition produce the correct answer, verified against ground truth computed
independently from the DisMech YAML files?) and \emph{heuristic score} (a
partial-credit measure based on string matching, reported as a secondary
diagnostic of response quality).

\paragraph{Binary accuracy.}
The structured condition produced verified correct answers for all ten questions
(10/10). The RAG condition produced zero correct answers (0/10). For Category~3
aggregation questions (Q1--Q3), RAG either returned an explicit refusal or an
incomplete partial list of diseases from retrieved chunks; neither constitutes a
correct corpus-wide count. For Category~5 questions (Q4--Q6), RAG explicitly
declined all three, correctly identifying that the retrieved context provided
insufficient information to reason about absence. For Category~6 questions
(Q7--Q10), RAG produced either explicit refusals or rankings derived from the
retrieved subset alone, yielding wrong answers.

\paragraph{Heuristic scores.}
The partial-credit scorer captures gradations within the RAG condition---some
responses named a few correct diseases; others refused entirely---but
substantially \emph{underestimates} structured performance due to three
implementation artifacts: (1)~the number extractor picks up incidental counts
(e.g., the corpus total of 605~diseases) before the specific count in scope,
awarding partial credit instead of exact credit; (2)~disease-name matching
fails on YAML underscore conventions (\texttt{Fanconi\_Anemia} in the source
files vs.\ ``Fanconi Anemia'' in Claude's prose); and (3)~the category regex
cannot match \texttt{(none/unclassified)} against the literal \texttt{(none)}
token from the ground truth. After correcting for these artifacts, heuristic
structured scores converge to~1.0 across all questions. The reported heuristic
structured scores should therefore be read as lower bounds.

\begin{table}[t]
\centering
\caption{Accuracy by question under RAG and structured conditions.
Binary accuracy is the primary measure (0~=~incorrect or refused;
1~=~verified correct against ground truth). Heuristic score is a
partial-credit string-matching measure. Dagger~($\dagger$) marks
structured scores that are underestimates due to scorer artifacts
described in Section~\ref{sec:results-overall}; corrected values
are~1.0 in all cases.}
\label{tab:accuracy}
\small
\begin{tabular}{clcccc}
\toprule
 & & \multicolumn{2}{c}{\textbf{Binary accuracy}} &
   \multicolumn{2}{c}{\textbf{Heuristic score}} \\
\cmidrule(lr){3-4}\cmidrule(lr){5-6}
\textbf{ID} & \textbf{Category} & RAG & Structured & RAG & Structured \\
\midrule
Q1  & Aggregation    & 0 & \textbf{1} & 0.000 & \textbf{1.000} \\
Q2  & Aggregation    & 0 & \textbf{1} & 0.000 & \textbf{1.000} \\
Q3  & Aggregation    & 0 & \textbf{1} & 0.000 & \textbf{1.000} \\
\midrule
Q4  & Negative space & 0 & \textbf{1} & 0.024 & 0.500$^\dagger$ \\
Q5  & Negative space & 0 & \textbf{1} & 0.000 & \textbf{1.000} \\
Q6  & Negative space & 0 & \textbf{1} & 0.000 & 0.250$^\dagger$ \\
\midrule
Q7  & Ranking        & 0 & \textbf{1} & 0.140 & 0.560$^\dagger$ \\
Q8  & Ranking        & 0 & \textbf{1} & 0.000 & 0.000$^\dagger$ \\
Q9  & Ranking        & 0 & \textbf{1} & 0.140 & 0.560$^\dagger$ \\
Q10 & Ranking        & 0 & \textbf{1} & 0.000 & 0.472$^\dagger$ \\
\midrule
\textbf{Mean} & & \textbf{0/10} & \textbf{10/10} & 0.026 & 0.634$^\dagger$ \\
\bottomrule
\end{tabular}
\end{table}

\subsection{Qualitative Analysis of RAG Failure Modes}
\label{sec:results-qualitative}

The RAG condition failed not by returning plausible-but-wrong numbers, but by
correctly recognising architectural impossibility. Claude's responses under RAG
fell into three distinct patterns.

\paragraph{Pattern~1---Explicit refusal (6/10 questions).}
For Q2, Q4, Q5, Q6, Q8, and Q9, the model identified that the retrieved context
was insufficient and declined to answer:

\begin{quote}
\emph{``Based on the provided context, I cannot determine the answer to this
question. The retrieved passages contain only PMIDs, abstract references, and
one disease description title, but none of them provide information about
disease categorisation as `Mendelian' or counts of diseases by
category.''}\quad---Q2 (Mendelian count; correct answer: 186)
\end{quote}

\begin{quote}
\emph{``I cannot answer this question. The retrieved passages only show abstract
entries with PMIDs and a few mechanism descriptions\ldots they do not provide a
comprehensive view of the citation counts across all sections of
DisMech.''}\quad---Q9 (top-5 by PMID citations; correct answer: Fanconi Anemia
with 254~citations)
\end{quote}

These refusals are epistemically correct: the retriever cannot surface what is
not there. They are nonetheless scored as failures because the benchmark
evaluates answer correctness, not calibration.

\paragraph{Pattern~2---Partial enumeration with wrong count (3/10 questions).}
For Q1, Q3, and Q7, retrieval surfaced a semantically relevant subset and the
model enumerated what it found. However, the retrieved sample is not the full
corpus. Q1 returned 3--4 TGF-$\beta$ diseases from retrieved chunks; the
correct answer is~26. Q7 identified Diabetes mellitus as having ``at least~2
mechanisms'' from the retrieved passages; the correct answer is~32:

\begin{quote}
\emph{``Based on the retrieved context, I can only identify a few diseases and
their individual mechanisms: Diabetes mellitus---at least 2~mechanisms
mentioned\ldots''}\quad---Q7 (correct answer: Diabetes mellitus, 32~mechanisms)
\end{quote}

The model correctly hedged with ``at least,'' signalling awareness that
retrieval is incomplete. This is the canonical failure mode for aggregation
queries under RAG: the retrieved sample provides a lower bound, not the answer.

\paragraph{Pattern~3---Vocabulary confusion (1/10 questions).}
For Q10 (full category distribution across 36~categories), retrieval returned
passages in which disease names incidentally resembled category labels:

\begin{quote}
\emph{``Based on the retrieved context, I can only identify two disease
categories that are explicitly mentioned: (1)~Infectious Disease,
(2)~Beh\c{c}et's Disease.''}\quad---Q10 (correct answer: 36~categories,
including Mendelian with 186~diseases)
\end{quote}

Beh\c{c}et's Disease is a disease name, not a category. The retriever surfaced
it from a passage that used the word in a context that the model interpreted
as categorical, demonstrating how semantic similarity can introduce category
errors when the schema distinction between entities and their classifications
is not structurally represented.

\subsection{Enriched-Tier Questions (Q11--Q13)}
\label{sec:results-enriched}

The three enriched-tier questions probe the nine-tier DisMech knowledge graph
extended beyond pathophysiology mechanisms to cover phenotypes with HPO
identifiers, genetic entries, treatments, evidence items, and causal edges.
To give RAG its best possible chance, the Qdrant index was extended from
11{,}456 to 20{,}104 vectors by adding 4{,}629 phenotype-description texts,
2{,}437 treatment-description texts, and 1{,}582 genetic-description texts.
Structured queries for Q11--Q13 use TypeDB graph traversal directly, replacing
the YAML-scan approach used for Q1--Q10.

\paragraph{Q11 (Category~3 --- gene-count aggregation).}
\textit{``How many diseases in DisMech have FGFR3 explicitly listed as a
causative gene in the structured gene annotations of their pathophysiology
mechanisms?''}

Structured (TypeDB 3-hop traversal:
$\mathit{disease} \to \mathit{pathophysiology\text{-}rel} \to
\mathit{pathophysiology} \to \mathit{gene} \to
\mathit{genedescriptor}[\texttt{preferred\text{-}term}=\text{``FGFR3''}]$):
\textbf{score~1.000}. The query returned the exact set of five diseases:
Achondroplasia, Hypochondroplasia, SADDAN, Thanatophoric Dysplasia Type~1, and
Thanatophoric Dysplasia Type~2.

RAG (top-20 from 20{,}104 vectors): \textbf{score~0.000}. The model identified
three diseases from mechanism description prose but missed Hypochondroplasia,
misidentified Diastrophic Dysplasia (which mentions FGFR3 in mechanism text
but not in the structured gene annotation field), and concluded ``at least
3~diseases''---a 40\% undercount with an unknown bound.

\paragraph{Q12 (Category~5 --- cross-tier absence detection).}
\textit{``How many diseases in DisMech have at least one HPO phenotype term
documented but NO documented genetic entries?''}

Structured (TypeDB set subtraction across two enriched tiers:
$\mathit{has\_HPO\_phenotype} \setminus \mathit{has\_genetic\_entry}$):
\textbf{score~1.000}. Query returned count~116. The TypeDB query expresses this
two-tier negation in 6~lines of declarative TypeQL; the equivalent YAML scan
requires approximately 20~lines of nested dictionary traversal.

RAG: \textbf{score~0.000}. The model correctly stated that it could not answer
from the retrieved context, explaining that the passages did not contain
sufficient information about co-occurrence of HPO annotations and absence of
genetic entries. This is the expected failure mode: no text in the corpus
represents the absence of a genetic entry.

\paragraph{Q13 (Category~6 --- phenotype-count ranking).}
\textit{``Which 5 diseases in DisMech have the most documented HPO phenotypes?
Give the count for each.''}

Structured (TypeDB 3-hop aggregation with \texttt{ORDER~BY}/\texttt{LIMIT~5}):
\textbf{score~0.700}. The top-5 ranking was: Fanconi Anemia (90~HPO
phenotypes), CLOVES Syndrome (38), Wilson Disease (28), Glycogen Storage
Disease Type~I (27), Cystic Fibrosis (25).

RAG: \textbf{score~0.000}. The model correctly declined: ``I cannot answer
this question. The retrieved passages do not contain information about HPO
phenotype counts.'' Even with 4{,}629 phenotype-description texts in the index,
retrieval returns semantically relevant phenotype entries---not the diseases
with the highest phenotype cardinality.

\begin{table}[h]
\caption{Q11--Q13 enriched-tier scores (one run per condition).}
\label{tab:enriched}
\small
\centering
\begin{tabular}{llccc}
\toprule
QID & Category & GT Answer & RAG & Structured \\
\midrule
Q11 & aggregation & 5 diseases   & 0.000 & 1.000 \\
Q12 & absence     & 116 diseases & 0.000 & 1.000 \\
Q13 & ranking     & top-5 by HPO count & 0.000 & 0.700 \\
\midrule
\multicolumn{3}{l}{\textit{Mean}} & 0.000 & 0.900 \\
\bottomrule
\end{tabular}
\end{table}

The enriched-tier questions confirm that the structural barriers identified for
Q1--Q10 extend uniformly to deeper ontological tiers. TypeDB's graph traversal
expresses 3-hop queries and cross-tier set subtraction as first-class
declarative operations; RAG cannot perform any of them regardless of corpus
size. The updated 20{,}104-vector index was the largest corpus RAG could
receive---and it produced zero correct answers.

\subsection{Why Retrieval Cannot Close the Gap}
\label{sec:results-analysis}

The RAG failures are not attributable to retrieval quality. The Qdrant
collection contained 11{,}456 vectors covering all mechanism description texts,
all disease summaries, and all 8{,}217 PubMed abstracts cited in DisMech,
embedded with \texttt{voyage-4-large} (1024~dimensions). Every piece of
information in the knowledge base was represented in the retrieval index.

For Category~3 (aggregation), a perfect retriever returning all 186~Mendelian
disease entries would still require Claude to count 186~prose passages correctly
under a single context window---with no guarantee that the 187th is not missing.
Increasing $k$ asymptotically approaches a full corpus scan but cannot verify
completeness; it also exceeds practical context limits long before reaching
corpus-wide coverage.

For Category~5 (absence detection), the problem is definitional: an absent
field generates no text. A disease with no documented treatments produces no
treatment-description vector. No retrieval strategy can return evidence of
non-existence from a corpus of positive assertions. This is the
\textsc{not~exists} barrier identified by~\citet{vardi1982complexity}: negation
is not expressible in the positive relational fragment corresponding to
retrieval-and-rank.

For Category~6 (global ranking), retrieval is biased by semantic relevance
rather than the structural property being ranked. Diseases with richer PubMed
literature are retrieved more frequently regardless of their mechanism counts,
introducing systematic bias into any ranking derived from retrieved chunks.
A disease with 32~documented mechanisms but sparse abstract coverage
(Folliculitis, ranked 4th) is outcompeted in retrieval by diseases with
extensive literature but fewer mechanisms.

These failure modes are independent of embedding model quality, chunk size, or
retrieval depth. They arise because aggregation, negation, and global ordering
are relational algebra primitives~\citep{codd1970relational,vardi1982complexity}
that have no retrieval-based equivalents.

\subsection{Cost and Practicality}
\label{sec:results-cost}

The full benchmark (10~questions $\times$ 2~conditions, one run per condition)
consumed 10{,}198 input tokens and 4{,}045 output tokens via the Claude Sonnet
4.6 API, at a total cost of \$0.09. The structured condition required no vector
index infrastructure and completed in under two minutes on a standard laptop.
The RAG condition required building an 11{,}456-vector Qdrant index
(approximately five minutes at Voyage~AI batch rates) plus one embedding and
retrieval call per question. For production deployments queried repeatedly, the
structured approach is also strictly cheaper per query once the knowledge graph
is loaded.
